\honest{Einstein's Speed}{Eliezer Yudkowsky}{2008}

Yesterday I argued that scientists must use their own judgment to decide which ideas to test. The easy case is a regularity spotted in old data. The hard case is Einstein, who, I say, ``had almost no experimental data to go on, except the precession of Mercury's perihelion.''

He got special relativity from ``Mach's Principle'': you should not be able to tell your speed inside a closed room. \nb{``Mach's principle'' is Einstein's 1918 name for a later idea, that the distribution of matter determines inertia. The principle for uniform motion is the principle of relativity.} Then he set out to do the same for acceleration. Since inertial and gravitational mass are always equal, gravity should be a kind of inertia, in curved spacetime. \nb{That equality is a measured fact about gravity, which Eötvös had measured in 1889, and Einstein built on it as a fact of experience.} He wrote down the simplest theory with the properties he thought it ought to have, and Mercury's orbit came out. \nb{Asked in 1993 whether general relativity is ``perfectly Machian,'' 22 of 25 experts said no. So one guiding property is not fully in the finished theory.}

Einstein, I say, ``didn't even use the perihelion precession of Mercury, except for verification,'' and made no mistakes ``that would have required further experimental evidence to pull him back on track.'' \nb{In 1913 Einstein computed Mercury for his earlier Entwurf theory and got too little, and in 1915 he named that failure as one of three reasons for giving the theory up. The two-year detour was corrected without new observations, as the post says.}

In Bayesian terms, Einstein learned the ``character of physical law'' from other laws and used it to predict a new one. \nb{Einstein described his method in much the same terms in 1919, and called its principles ``empirically discovered.''}

Asked about the eclipse, Einstein said, ``Then I would feel sorry for the good Lord. The theory is correct.'' As I argued in ``Einstein's Arrogance,'' he must have had the evidence. So: ``Wait for an occasion where they are wrong!'' \nb{There was one. In 1914 Einstein wrote of the Entwurf theory that he no longer doubted its correctness, ``regardless of whether the observation of the solar eclipse will be successful or not.'' It was wrong.} Einstein told Bohr that God does not play dice, and Bohr told Einstein not to tell God what to do. \nb{That is the popular version; Bohr's line is disputed.}

Science made the theory wait for the eclipse, and that measurement ``had favored GR essentially by luck.'' \nb{Contested. A 1979 remeasurement of the plates supported the original conclusion that the Newtonian value was ruled out. And on Stephen Brush's reading, physicists gave the eclipse no more weight than Mercury.}

History remembers whoever got there first, so do not think armchair reasoning easy or reliable. Use every scrap of evidence. For careful reasoning, study evolutionary psychology, whose practitioners, perhaps, ``do later find out if they were right or wrong.'' \nb{The post gives no example, and whether the field's hypotheses get tested is its main point of dispute.} At one in a million, there are ``at least six thousand potential Einsteins running around today.''
